\section{Applying the protocol to external projects}
\label{app:external-projects}

Projects funded by external organisations (e.g., NWO, Horizon Europe, industry, or public-private 1-on-1 
partnerships) are subject to contractual obligations defined by the funding agreement, collaboration agreement, or 
employment contract. Where these obligations differ from the guidance provided in this protocol, the contractual 
requirements take precedence.

This appendix explains how the Project Management Protocol applies to externally funded projects and highlights the 
main adaptations required to comply with funding conditions while remaining consistent with the eScience Center's 
internal project management practices. Further information on the objectives and scope of external projects and the 
approval process for pursuing externally funded projects is available in the eScience Center's External Projects 
Policy~\cite{nlesc-external-projects-policy}.

\subsection{Project initiation}
The starting point of an externally funded project is determined by the contractual arrangements with the funding 
organisation (e.g. consortium agreement or 1-on-1 contract). Depending on the funding scheme, project initiation 
may consist of a bilateral kick-off meeting with the LA, or of a formal consortium kick-off meeting involving all 
project partners. Regardless of the external arrangements, the eScience Center follows the initiation activities 
described in this protocol, unless they conflict with contractual requirements.

Some deviations may include: 
\vspace{-0.2cm}
\begin{itemize}
 \item The project start date is defined by the funding agreement.
 \item Consortium kick-off meetings may replace or complement the project kick-off described in this protocol.
 \item Additional project documentation (e.g., Grant Agreement, Consortium Agreement, IP Agreement) should be 
 registered together with the project records in the corresponding project portfolio folder.
 \item A separate Technology Plan may not be required if the project proposal or contract arrangement already defines 
 the project's technical choices and implementation approach. 
\end{itemize}
\vspace{-0.2cm}

\subsection{Project execution}
Unless specified otherwise by the funding agreement or consortium arrangements, the project execution practices 
described in this protocol remain applicable. This includes regular project and internal coordination meetings, 
maintenance of the project log and RSD page, recording of project hours, and the software quality and sustainability 
practices promoted by the eScience Center.

In addition to the above, some adaptations may be required to comply with the contractual obligations of the funding 
organisation. These may include:
\vspace{-0.2cm}
\begin{itemize}
 \item Project reporting, deliverables, milestones, and review processes are governed by the funding agreement and 
 take precedence over the reporting recommendations described in this protocol. These activities are usually handled 
 and led by the project coordinator of the funding organization. 
 \item Financial information required for external reports should be provided by Finance. Lead RSE should keep key 
 project and reporting documentation records in the project portfolio folder.
 \item DMPs, SMPs and other project management artefacts should follow the format and schedule required by the 
 funding organisation.
 \item Project activities such as dissemination, communication, community building, and similar activities should be 
 carried out in accordance with the contractual requirements.
\end{itemize}
\vspace{-0.2cm}

\subsection{Project closing}
The project closing activities described in this protocol remain applicable. In addition to the above, the following 
adaptations may be required to comply with the contractual obligations of the funding organisation:

\vspace{-0.2cm}
\begin{itemize}
  \item The project closing meeting may be organised by the funding organisation or project coordinator and may replace 
  or complement the closing meeting described in this protocol.
  \item The final project report is typically coordinated by the funding organisation or project coordinator. The Lead 
  RSE should retain a copy of the approved report and, together with Finance, complete the internal administrative 
  procedures required to close the project.
  \item 1-on-1 service contracts may not require a formal end report. In such cases, Lead RSE ensures the project 
  closure follows the contractual ending procedure agreed with the partner. 
\end{itemize}
\vspace{-0.2cm}